\begin{longtable}{p{3.0cm}p{5.4cm}p{4.8cm}}
\caption{Calidad de los registros y tratamiento previsto}\label{tab:a3}\\
\toprule
Aspecto & Evidencia & Tratamiento previsto \\
\midrule
\endfirsthead
\multicolumn{3}{l}{\footnotesize\itshape (continuación)}\\
\toprule
Aspecto & Evidencia & Tratamiento previsto \\
\midrule
\endhead
\bottomrule
\endfoot
Variables del evento sin vacíos & Fecha, hora, provincia, cantón, zona, tipo de siniestro, causa y ente de control & Se usan tal como están \\
Jerarquía y ubicación de la vía & \texttt{ND} en 11.445 registros (66,4\,\%); 14 con valor 0 y 1 vacío & Se trata como categoría propia; el valor 0 se agrupa con \texttt{ND} \\
Tipo del primer vehículo no identificado & 3.870 registros (22,5\,\%) & Se trata como categoría propia y se conserva \\
Tipo de siniestro con dos escrituras & 10 registros con \texttt{colisión} en minúscula & Se unifican con \texttt{COLISIÓN} \\
Hora almacenada como texto & 5 registros & Se convierten a hora y se incorporan a la franja horaria \\
Identificador repetido & 1 identificador (\texttt{MCU00139082026}) en dos siniestros de fechas distintas & Se conservan ambas filas; no se usa el identificador como llave \\
Filas duplicadas & 0 & No se elimina ninguna fila \\
Coordenadas & Sin vacíos y dentro de un rango plausible & No se usan (restricción R-07) \\
\end{longtable}
